%% template_data_integration.tex
%% EuMINe DataBridge Hackathon 2026 — Data Integration Report Template
%%
%% Fill in every TODO section.
%% Compile with: pdflatex template_data_integration.tex
%%               (run twice for cross-references)

\documentclass[10pt, a4paper]{article}
\usepackage{eumine_hackathon}

%% ── Title block ────────────────────────────────────────────────────────────
\title{%
  \textbf{Data Integration Report}\\[4pt]
  {\large EuMINe DataBridge Hackathon \hackathonyear}
}
\author{
  TODO: Team Name \\
  \small TODO: Institution(s) \\
  \small TODO: Contact email
}
\date{TODO: Date submitted}

%% ───────────────────────────────────────────────────────────────────────────
\begin{document}
\maketitle
\thispagestyle{fancy}

% ── Abstract ──────────────────────────────────────────────────────────────
\begin{abstract}
TODO: One paragraph summarising the data sources used, how they were
harmonised, and the key characteristics of the final training dataset.
\end{abstract}

\tableofcontents
\vspace{1em}

% ── 1. Data Sources ────────────────────────────────────────────────────────
\section{Data Sources Used}

\begin{table}[h!]
\centering
\caption{Summary of data sources included in the training dataset.}
\label{tab:sources}
\begin{tabular}{@{}lllll@{}}
\toprule
Source & Version / date & Entries used & Properties & DFT functional \\
\midrule
Materials Project  & TODO & TODO & \(E_f\), \(E_g\) & PBE (GGA) \\
JARVIS-DFT dft\_3d & TODO & TODO & \(E_f\), \(E_g\) & OptB88vdW \\
TODO (other)       & TODO & TODO & TODO             & TODO \\
\bottomrule
\end{tabular}
\end{table}

\subsection{Materials Project}
TODO: Describe what fields were downloaded, filters applied (e.g.\
\texttt{has\_props=["thermo","electronic\_structure"]}), and any known
caveats (band gaps from HSE vs.\ GGA, deprecated entries, etc.).

\subsection{JARVIS-DFT}
TODO: Describe the \texttt{dft\_3d} dataset version used, key fields
(\texttt{formation\_energy\_peratom}, \texttt{optb88vdw\_bandgap}), and
notable differences from MP (functional, pseudopotentials, k-point density).

\subsection{Additional Sources (if any)}
TODO: Describe any extra databases used.

% ── 2. Harmonisation Strategy ─────────────────────────────────────────────
\section{Harmonisation Strategy}

\subsection{Structure Matching}
TODO: Explain matching logic (e.g.\ reduced formula + space group number, plus
\(\pm 2\) site tolerance). Report number of raw/filtered/final matches.

\subsection{Handling Duplicates}
TODO: Describe deduplication (e.g.\ keep lowest-energy polymorph per
formula+spacegroup combination, or random selection with fixed seed).

\subsection{Label Assignment}
TODO: Explain which source is used as the primary label for training (e.g.\
MP value primary, JARVIS as fallback if MP value is missing). Justify the
choice.

\subsection{Methodological Bias Correction (optional)}
TODO: If you applied any bias correction (e.g.\ a linear offset between MP
and JARVIS formation energies), describe the method and its effect on
property distributions.

% ── 3. Data Card ──────────────────────────────────────────────────────────
\section{Data Card}

\begin{table}[h!]
\centering
\caption{Dataset statistics (training split).}
\label{tab:datacard}
\begin{tabular}{@{}ll@{}}
\toprule
Property & Value \\
\midrule
Total entries (train / val / test) & TODO / TODO / TODO \\
Unique chemical systems            & TODO \\
Space group count                  & TODO \\
Median \(|\Delta E_f|\) MP\,--\,JARVIS & TODO eV/atom \\
Median \(|\Delta E_g|\) MP\,--\,JARVIS & TODO eV \\
Splitting strategy                 & TODO (stratified / random) \\
Stratification variable            & TODO \\
\bottomrule
\end{tabular}
\end{table}

TODO: You may also include histograms of formation energy and band gap
(training set) generated from \texttt{bridge\_dataset\_train.csv}.

% ── 4. Limitations and Failure Cases ─────────────────────────────────────
\section{Limitations and Failure Cases}

\begin{itemize}
  \item TODO: Compositional coverage gaps (e.g.\ rare-earth elements
    under-represented because JARVIS coverage is low).
  \item TODO: Known issues with band gap labels (semiconductors with
    mis-classified gap due to PBE underestimation).
  \item TODO: Structural features that may cause featuriser failures
    (e.g.\ disordered structures, very large supercells).
  \item TODO: Any materials excluded and why.
\end{itemize}

% ── References ────────────────────────────────────────────────────────────
\section*{References}

\begin{enumerate}
  \item Jain et al., \textit{APL Mater.} \textbf{1}, 011002 (2013) —
    Materials Project.
  \item Choudhary et al., \textit{npj Comput. Mater.} \textbf{6}, 173 (2020)
    — JARVIS-DFT.
  \item TODO: Add any additional references.
\end{enumerate}

\end{document}
